\par\medskip
\begin{center}
\textls[90]{c) Die Jahre 1809, 1810. Der bilineare Algorithmus. Persönliches.}
\end{center}
\noindent\hspace*{\parindent}Aus dem Juni 1809 haben wir zwei auf die Theorie der elliptischen
Funktionen bezügliche Tagebuchnotizen: 139. vom 20. Juni, »Die zu den
agMitteln gehörigen Reihen weiter entwickelt« und 140. vom 29. Juni, »Die
Fünfteilung für die agMittel erledigt«.
\par
Die Werke III, S. 446—460 abgedruckte Abhandlung stammt aus dem
im Oktober 1805 begonnenen Handbuche 18, Bd, wo sie S. 221—233 unmittelbar
nach einer astronomischen Rechnung aufgezeichnet ist, der die Bemerkung
beigefügt ist: »geendiget d. 28. April 1809«. Da sie mit den Worten
beginnt: »Die Theoreme in Beziehung auf diejenigen Reihen und unendlichen
Produkte, welche zu der Theorie des arithmetisch-geometrischen Mittels gehören,
ordnen wir so«, ist ihr Zusammenhang mit der Tagebuchnotiz 139.
unverkennbar, ferner beweist die Behandlung der Fünfteilung der Perioden
auf S. 456—460, dass die Schlusspartie dieser Abhandlung um die Zeit der
Tagebuchaufzeichnung 140. entstanden ist. Dem Inhalte nach mit dieser Abhandlung
verwandt ist die den Seiten 37—40 der Scheda An entnommene,
Werke X, 1, S. 213—216 abgedruckte Aufzeichnung, die dadurch, dass sich
auf S. 35 der Scheda An am Schluss einer astronomischen Rechnung die Bemerkung\par
\HistoricalFoot{X 2 Abh. 2.}{17}
\clearpage\SourceLeaf{215}\ArithmeticHead{130}{L}{L. SCHLESINGER, ÜBER GAUSS' FUNKTIONENTHEORETISCHE ARBEITEN.}
\noindent findet: »geendigt d. 2. Mai 1809« mit Sicherheit auf die Zeit der
Tagebuchnotizen 139., 140. zu datieren ist.
\par
Neben uns bereits bekannten Formeln zur Theorie des agM. und den
Umformungen von Reihen in Produkte enthält die Aufzeichnung der Scheda
An als neuartiges Element einen Algorithmus, der eine Erweiterung des agM.
in dem Sinne darstellt, dass er sich in ähnlicher Weise auf die Thetafunktionen
von zwei Veränderlichen aufbaut, wie das agM. auf die summatorischen
Funktionen \(p,q,r\). Dieser seiner Form nach \textls[100]{bilineare} Algorithmus ergibt
sich folgendermassen: \textsc{Gauss} setzt (Werke X, 1, S. 214, Gl. (5))\footnote{Wir schreiben die Formeln in den gewohnten Zeichen; bei \textsc{Gauss} steht \(x^4\) an Stelle unseres \(x\) und \(y=e^{\pi i\psi}\) statt unseres \(\psi\).}
\begingroup\makeatletter\tagsleft@true\let\veqno\leqno\makeatother
\begin{equation*}
A=H\vartheta_{00}(\psi|x)^2,\quad B=H\vartheta_{01}(\psi|x)^2.\tag{1}
\end{equation*}\endgroup
\noindent Ähnlich wie das agM. durch die Folge
\[
\frac{b}{a}=\frac{q(x)^2}{p(x)^2},\quad
\frac{b_1}{a_1}=\frac{q(x^2)^2}{p(x^2)^2},\quad
\frac{b_2}{a_2}=\frac{q(x^4)^2}{p(x^4)^2},\ \ldots
\]
mit \(a_1=\tfrac12(a+b),\ a_2=\tfrac12(a_1+b_1),\ \ldots\) definiert ist, bildet \textsc{Gauss} nun gemäss
(1) die Folge
\begingroup\makeatletter\tagsleft@true\let\veqno\leqno\makeatother
\begin{equation*}
\frac{B}{A}=\frac{\vartheta_{01}(\psi|x)^2}{\vartheta_{00}(\psi|x)^2},\quad
\frac{B_1}{A_1}=\frac{\vartheta_{01}(2\psi|x^2)^2}{\vartheta_{00}(2\psi|x^2)^2},\quad
\frac{B_2}{A_2}=\frac{\vartheta_{01}(4\psi|x^4)^2}{\vartheta_{00}(4\psi|x^4)^2},\ \ldots\tag{2}
\end{equation*}\endgroup
\noindent und setzt zur Sicherung einer kräftigen Konvergenz wie beim agM.
\begingroup\makeatletter\tagsleft@true\let\veqno\leqno\makeatother
\begin{equation*}
A_1=\tfrac12(A+B),\quad A_2=\tfrac12(A_1+B_1),\ \ldots.\tag{3}
\end{equation*}\endgroup
\noindent Mit den Formeln der \textsc{Landenschen} Transformation, wie sie in der Scheda Ac
für die Thetareihen gegeben wurden (siehe unsere Formeln (24), (25), (25a)
Abschnitt IV, S. 78), folgt nunmehr aus (2)
\begingroup\makeatletter\tagsleft@true\let\veqno\leqno\makeatother
\begin{equation*}
\frac{B_1}{A_1}=\frac{AB}{\left(\frac{A+B}{2}\right)^2}\frac{a_1}{b_1},\quad
\frac{B_2}{A_2}=\frac{A_1B_1}{\left(\frac{A_1+B_1}{2}\right)^2}\frac{a_2}{b_2},\ \ldots,\tag{4}
\end{equation*}\endgroup
\noindent also mit Rücksicht auf die Festsetzung (3)
\begingroup\makeatletter\tagsleft@true\let\veqno\leqno\makeatother
\begin{equation*}
B_1=\frac{2AB}{A+B}\frac{a_1}{b_1},\quad
B_2=\frac{2A_1B_1}{A_1+B_1}\frac{a_2}{b_2},\ \ldots,\tag{5}
\end{equation*}\endgroup
\noindent was in Verbindung mit (3) den von \textsc{Gauss} a. a. O. S. 214, Gl. (3) angegebenen
bilinearen Algorithmus darstellt. Ähnlich wie man aus dem Ansatz
\[
a=\mu p(x)^2,\quad b=\mu q(x)^2
\]
für das agM. bei willkürlich gegebenen \(a,b\) die zugehörigen Grössen \(\mu,x\)\par
\clearpage\SourceLeaf{216}\ArithmeticHead{131}{R}{VIc. DER BILINEARE ALGORITHMUS.}
\noindent bestimmen kann, nämlich (vergl. oben S. 107)
\[
\mu=M(a,b),\quad x=\frac{a^2-b^2}{16a_1a_2}
\left(\frac{a_2}{a_3}\right)^{\frac12}
\left(\frac{a_3}{a_4}\right)^{\frac14}\ldots
=e^{-\pi\frac{M(a,b)}{M(a,c)}},
\]
so zeigt \textsc{Gauss}, dass man aus dem Ansatz (1) die Grössen \(H\) und \(\psi\) zu berechnen
imstande ist. Die erstere Grösse ergibt sich unmittelbar\footnote{Vergl. die Darstellung bei \textsc{Geppert} a. a. O. \textsuperscript{83)}, § 42, S. 177 ff.}, wenn
man aus den gemäss (2) und (3) folgenden Ausdrücken
\[
A_1=H\sqrt{\frac{a_1}{\mu}}\vartheta_{00}(2\psi|x^2),\quad
A_2=H\sqrt{\frac{a_1}{\mu}\frac{a_2}{\mu}}\,
\frac{\vartheta_{00}(4\psi|x^4)}{\vartheta_{00}(2\psi|x^2)},\ \ldots
\]
das Produkt \(A_1^{\frac12}A_2^{\frac14}\ldots\) bildet, in der Form (Werke X, 1, S. 214, (4))
\begingroup\makeatletter\tagsleft@true\let\veqno\leqno\makeatother
\begin{equation*}
H=\mu\frac{A_1^{\frac12}A_2^{\frac14}\ldots}{a_1^{\frac12}a_2^{\frac14}\ldots}.\tag{6}
\end{equation*}\endgroup
\noindent Um ferner \(\psi\) oder \(y\) zu erhalten, stellt \textsc{Gauss} einen trigonometrischen, für das
logarithmische Rechnen sehr geeigneten Algorithmus auf (a. a. O. Art. 2),
nämlich
\begingroup\makeatletter\tagsleft@true\let\veqno\leqno\makeatother
\begin{equation*}
\left\{\begin{aligned}
&\frac{b}{a}=\cos 2M,\quad\frac{b_1}{a_1}=\cos 2M_1,\quad\frac{b_2}{a_2}=\cos 2M_2,\ \ldots\\
&\sin\Phi\sqrt{\sin 2M}=\operatorname{tg}\Psi,\quad
\sin 2\Psi_1=\operatorname{tg}\Psi\sqrt{\sin 2M},\\
&\sin\Phi_1\sqrt{\sin 2M_1}=\operatorname{tg}\Psi_1,\quad
\sin 2\Psi_2=\operatorname{tg}\Psi_1\sqrt{\sin 2M_1}\\
&\sin\Phi_2\sqrt{\sin 2M_2}=\operatorname{tg}\Psi_2,\\
&\hspace{10em}\text{usw.}
\end{aligned}\right.\tag{7}
\end{equation*}\endgroup
\noindent der, wenn man (a. a. O. S. 215)
\begingroup\makeatletter\tagsleft@true\let\veqno\leqno\makeatother
\begin{equation*}
\operatorname{tg}\Psi=\frac{\vartheta_{10}(\psi|x)}{\vartheta_{00}(\psi|x)}\tag{8}
\end{equation*}\endgroup
\noindent setzt, darauf hinauskommt, dass
\[
\begin{aligned}
\sin\Phi&=\frac{p(x)}{r(x)}\frac{\vartheta_{10}(\psi|x)}{\vartheta_{00}(\psi|x)}
=\operatorname{Cl}\varphi,\quad\varphi=\pi\psi q(x)^2,\\
\sin\Phi_1&=\frac{p(x^2)}{r(x^2)}\frac{\vartheta_{10}(\psi|x^2)}{\vartheta_{00}(\psi|x^2)}
=\operatorname{Cl}\varphi_1,\quad\varphi_1=\pi\psi q(x^2)^2
\end{aligned}
\]
ist, also (vergl. oben S. 109) auf die \textsc{Gausssche} Transformation
\[
\sin\Phi=\frac{2a\sin\Phi_1}{a+b+(a-b)\sin^2\Phi_1}.
\]
\HistoricalFoot{}{17\textsuperscript{*}}
\clearpage\SourceLeaf{217}\ArithmeticHead{132}{L}{L. SCHLESINGER, ÜBER GAUSS' FUNKTIONENTHEORETISCHE ARBEITEN.}
\noindent Durch wiederholte Anwendung erhält man
\[
\sin\Phi_n=\operatorname{Cl}\varphi_n,\quad\varphi_n=\pi\psi q(x^{2^n})^2,
\]
so dass also
\begingroup\makeatletter\tagsleft@true\let\veqno\leqno\makeatother
\begin{equation*}
\lim_{n\to\infty}\sin\Phi_n=\cos\psi\pi\tag{9}
\end{equation*}\endgroup
\noindent wird, was \textsc{Gauss} (a. a. O. S. 215, (7)) in der Form
\[
y+\frac{1}{y}=2\sin\Phi_\infty
\]
schreibt\footnote{Der Algorithmus (3), (5) wird von \textsc{Schering}, Werke III, S. 389, wohl erwähnt, \textsc{Schering} hat auch erkannt, dass \textsc{Gauss} diesen Algorithmus mit gutem Vorbedacht gerade so bildet, um eine gute Konvergenz zu erzielen (vergl. die Bemerkungen S. 389, Zeile 8, 7 v. u. und S. 390, Zeile 7, 8); trotzdem hat \textsc{Schering} an Stelle des \textsc{Gaussschen} einen überhaupt nicht von \textsc{Gauss} herrührenden Algorithmus von vier Grössen \(\alpha,\beta,\gamma,\delta\) in Werke III, S. 387 ff. Artikel 18—25 entwickelt, der früher stets als der \textsc{Gausssche} gegolten hat. Die Beziehungen zwischen dem \textsc{Gaussschen} und dem \textsc{Scheringschen} Algorithmus sind ja sehr einfach, indem die \textsc{Scheringschen} \(\alpha,\beta,\gamma,\delta\) direkt den Quadraten der vier Thetafunktionen, \(\vartheta_{00}(\psi|x)^2,\) \(\vartheta_{01}(\psi|x)^2,\ \vartheta_{10}(\psi|x)^2,\ -\vartheta_{11}(\psi|x)^2,\) proportional sind, aber \textsc{Schering} erweckt durch seine Darstellung den Anschein, dass für \textsc{Gauss} dieser Algorithmus das prius gewesen, und dass er mit Hilfe desselben zu den Thetafunktionen vorgedrungen sei. Aus den Angaben P. \textsc{Günthers} (Gött. Nachr. 1894, S. 94 u. 104) geht hervor, dass auch \textsc{Weierstrass} sich durch die Darstellung \textsc{Scherings} hatte irreführen lassen. Wir wissen jetzt, dass \textsc{Gauss} schon 9 Jahre im Besitze der Thetafunktionen war, als er von den Transformationsformeln zweiter Ordnung dieser Funktionen ausgehend, wahrscheinlich nur zum Zwecke einer für numerische Rechnungen gut brauchbaren Entwicklung, den Algorithmus (3), (5) aufstellte.\par
Vergl. die von den Thetafunktionen ausgehende Darstellung des \textsc{Scheringschen} Algorithmus, die v. \textsc{Dávid} 1903 gegeben hat, Mathem. Naturw. Berichte aus Ungarn XXV, 1907, S. 153 ff. Die historischen Angaben dieser Arbeit sind natürlich nach dem obigen zu berichtigen.}. Die Beziehung zwischen diesem Algorithmus und dem zur Bestimmung
des Integrals erster Gattung dienenden (Werke X, 1, S. 227—229)
liegt auf der Hand.
\par
Wir sehen, dass \textsc{Gauss} in seinen auf die elliptischen Transzendenten bezüglichen
Arbeiten aus den Jahren 1808—1809 sich im wesentlichen in demselben
Gedankenkreise bewegt, wie in der Zeit der gewaltigen Produktivität,
in die uns das \textit{Tagebuch} und die Scheda Ac einen Einblick gewährt haben.
Es wird manches geglättet, vieles Neue kommt hinzu, aber es ist doch eine
Tätigkeit, die mehr sammelnd und ordnend die vorhandenen Schätze in gangbare
Münze umprägen, als in den Tiefen nach neuen Goldadern schürfen will.
Vielleicht soll schon das Material für das grosse Werk über die transzendenten
Funktionen bereit gestellt werden, nachdem in der \textit{Theoria motus} (1809)\par
\clearpage\SourceLeaf{218}\ArithmeticHead{133}{R}{VIc. PERSÖNLICHES 1809, 1810.}
\noindent einige Früchte der astronomischen Arbeit der letzten Jahre der Öffentlichkeit
übergeben worden sind.
\par
Diese Vermutung wird durch den Umstand bestätigt, dass sich in dem
Handbuch 18, Bd, schon auf S. 52 ff. (Werke X, 1, S. 326) und dann in unmittelbarem
Anschluss an die erwähnte, Werke III, S. 446—460 abgedruckte
Abhandlung auf den Seiten 234—235 (Werke X, 1, S. 332) Aufzeichnungen
über die Reihe \(F(\alpha,\beta,\gamma,x)\) finden, über die im folgenden Abschnitt berichtet
wird.
\par
In der zweiten Hälfte des Jahres 1809 und in dem ganzen folgenden
Jahre haben \textsc{Gauss}' Lebensschicksale hemmend auf seine wissenschaftliche
Produktion gewirkt. Die Verhältnisse in dem ephemeren Königreich Westfalen
mögen wenig erfreulich gewesen sein; der immer wieder versprochene
Neubau der Sternwarte kommt nicht zur Ausführung, und in einem Briefe an
\textsc{Olbers} vom 4. Oktober 1809\footnote{\textit{Briefwechsel Olbers-Gauss} I, S. 441.} spricht \textsc{Gauss} sogar die Befürchtung aus, dass
mit Ende des Jahres 1809 die Auflösung der Sozietät der Wissenschaften
bevorstehen könnte. Weitaus tiefer aber ging die Wirkung der Vorgänge in
\textsc{Gauss} engster Familie. Nach der Geburt eines Sohnes (10. September 1809)
konnte die Frau sich nicht wieder erholen und starb am 11. Oktober 1809.
»Gestern Abend«, schreibt er am 12. Oktober an \textsc{Olbers}\footnote{\textit{Briefwechsel Olbers-Gauss} I, S. 442; vergl. auch »\textsc{Gauss}' Totenklage um seine Frau« in H. \textsc{Mack}, \textit{C. F. Gauss und die Seinen}, Braunschweig 1927, S. 16, 17.}, »um 8 Uhr habe
ich ihr die Engelsaugen, in denen ich seit fünf Jahren einen Himmel fand,
zugedrückt. Der Himmel gebe mir Kraft, diesen Schlag zu ertragen«. Am
14. Dezember 1809 schreibt er an \textsc{Schumacher}\footnote{\textit{Briefwechsel Gauss-Schumacher} I, S. 17.}: »Mit dem Arbeiten will es
noch nicht recht bei mir gehen. Ich habe angefangen die biquadratischen
Reste wieder vorzunehmen, aber ich fühle mich noch immer zu wenig und
zu selten aufgelegt« und am 13. Januar 1810 an \textsc{Heyne}\footnote{Den Sekretär der Gesellschaft der Wissenschaften; Brief im Archiv der G. d. W.}: »Eine schon früher
angefangene Abhandlung hoffe ich noch diesen Winter der Sozietät übergeben
zu können, sobald teils meine nicht immer zu solchen Arbeiten aufgelegte
Gemütsstimmung, teils die zwei Kollegien, die ich in diesem Winter lese, es
erlauben«. Gemeint ist wohl die Abhandlung über \(F(\alpha,\beta,\gamma,x)\), die aber erst
zwei Jahre später und da nur in ihrem ersten Teile vorgelegt worden ist.\par
\clearpage\SourceLeaf{219}\ArithmeticHead{134}{L}{L. SCHLESINGER, ÜBER GAUSS' FUNKTIONENTHEORETISCHE ARBEITEN.}
\noindent Noch am 10. Februar 1810 schreibt er an \textsc{Schumacher}\footnote{\textit{Briefwechsel Gauss-Schumacher} I, S. 26.}: »Meine wissenschaftlichen
Beschäftigungen bedeuten in diesem Winter nicht viel. Meine
beiden Kollegia \mbox{.\ .\ .\ .}\ zerstückeln meine Zeit und machen mich für einen Teil
des Tages zum Arbeiten unlustig; es treibt mich dann immer, wenn ich gelesen
habe, aus dem Haus ins Weite.« Am 1. März 1810 starb das kaum
halbjährige Söhnchen, aber bald darauf verleiht ihm, wie er an \textsc{Olbers}
schreibt\footnote{Anfang April 1810, \textit{Briefwechsel Olbers-Gauss} I, S. 448.}, »eine sonderbare, fast romanhafte Concatenation von Umständen,
in denen ich fast Fingerzeige einer höheren Hand wahrzunehmen geneigt sein
möchte« neuen Lebensmut, und er geht am 4. August 1810 mit einer Freundin
seiner verstorbenen Frau, der Tochter seines Göttinger Kollegen \textsc{Waldeck},
eine zweite Ehe ein. Sehr allmählich kehrt auch die Arbeitskraft und Arbeitslust
wieder; am 25. November 1810 legt er der Gesellschaft der Wissenschaften
die grosse Abhandlung über die Pallas vor\footnote{Werke VI, S. 1 und \textit{Anzeige}, S. 61.}, und in dem Brief an
\textsc{Bessel} vom 21. Oktober 1810 (Werke X, 1, S. 360) gibt er eine Übersicht
über seine Methode, den Integrallogarithmus zu behandeln.
